% !PRINTED 659
\weberpage{659}

For the proof of the Principal Theorem I. everything now depends on
the proof of the two theorems that are to be inferred from the
properties of the functionals $\varphi$ and $\varepsilon$:

A. \gesperrt{The functionals $\varphi_n$ belong to the principal class.}

If we may assume this lemma, then $\varphi_n$ is associated with a
number $\alpha_n$, and we can represent the formula (14), with a
somewhat altered meaning of $\varepsilon_n$, in this
way:
\weq{15}{\omega_n = \varepsilon_n \alpha_n^m \prod^{p} (r^{m_2 n}, \eta)^m.}
In this formula, however, $\varepsilon_n$ is a \gesperrt{numerical unit}
of the field $\Omega_m$ which satisfies Theorem 9.

If, then, we can prove from Theorem 9. the second lemma as
well:

B. \gesperrt{A numerical unit $\varepsilon_n$ of the field
$\Omega_m$ which satisfies Theorem 9. is the $m$th
power of a unit in $\Omega_m$, multiplied by a power
of $r$;}

then we can extract the $m$th root in (15), and obtain, if by
$\varrho$ a root of unity of possibly higher degree than $m$,
and by $\alpha$ a number of the field $\Omega_m$, is denoted:
\weq{16}{(r, x_0) = \varrho \alpha \prod^{p} (r^{m_2}, \eta),}
where now nothing but cyclotomic numbers stand on the right, and
Theorem I. is therefore proved.

\begin{center}
\S.\,182.

\gesperrt{Proof of the first lemma for an odd $m$.}
\end{center}

In the preceding section we made the proof of Theorem I. depend
on two lemmas, whose proof is now still to be sought.

We formulate the first of these lemmas thus:

a) \gesperrt{If $\varphi_n$ is a system of conjugate
functionals of the field $\Omega_m$, different from zero, of
which it is known that}
\[
\varphi_n^m \text{ \gesperrt{and} } \varphi_n^{-1} \varphi_1^n
\]
\gesperrt{belong to the principal class, then the
functionals $\varphi$ themselves belong to the principal class.}
